\chapter{من الورق إلى كيان}
\label{ch:paper}

تتعلم كيف ينتقل عمل مؤسستك إلى كيان خطوة خطوة: ماذا تجمع أولاً، وما الذي تستخدمه جاهزاً، وما الذي تبنيه، وبأي ترتيب. ثم تطبّق ذلك على إجراء من عملك.

\section{التحول خطوة خطوة}
\label{sec:paper-steps}
\begin{rcslide}{التحول خطوة خطوة}
\sphase{المنهجية}
\slead{لا تنقل المؤسسة كلها دفعة واحدة. تنقل إجراءً واحداً كاملاً، ثم الذي يليه.}
\schain{الهيكل والناس ← السجلات والاستمارات ← الموافقات والإجراءات ← التقارير}
\smetric{نقطة البداية}{شراء برنامج ثم التكيف معه}{عملكم كما هو اليوم}
\smetric{الحجم}{المؤسسة كلها مرة واحدة}{إجراء واحد كامل}
\smetric{من يشارك}{فريق التقنية وحده}{من يعرف العمل نفسه}
\stake{كل إجراء تنقله كاملاً يعمل من يومه الأول، ويعلّمك ما تحتاجه للإجراء التالي.}
\end{rcslide}

\section{اجمع ما تعمل به اليوم}
\label{sec:paper-inventory}
\begin{rcslide}{اجمع ما تعمل به اليوم}
\sphase{المنهجية}
\slead{قبل أن تفتح كيان، اجمع الأوراق التي يدور عليها العمل. كل ورقة منها ستصبح شيئاً محدداً في كيان.}
\scard{building}{الهيكل التنظيمي}{الوحدات والمناصب ومن يشغلها.}
\scard{form}{الاستمارات}{كل استمارة يعبّئها موظف أو مراجع.}
\scard{sheet}{الدفاتر والجداول}{ما تتابعه المؤسسة: منشآت، مستفيدون، عقود، طلبات.}
\scard{flow}{مسار التوقيعات}{من يوقّع بعد من، ومتى يُعاد الطلب.}
\scard{chart}{التقارير الدورية}{الأرقام التي يطلبها المدير كل شهر.}
\stake{هذه الأوراق الخمس هي خريطة عملك. وكلما كانت أدق، كان البناء أسرع.}
\end{rcslide}

\section{لا تبنِ ما هو جاهز}
\label{sec:paper-ready}
\begin{rcslide}{لا تبنِ ما هو جاهز}
\sphase{المنهجية}
\slead{قبل أن تبني شيئاً، تأكد أنه غير موجود جاهزاً في كيان.}
\scompare{ما تعمل به اليوم}{الهيكل وملفات الموظفين \\ الإجازات والدوام \\ دفتر الوارد والصادر \\ مجلدات الملفات المشتركة \\ التكليفات بالرسائل}{التطبيق الجاهز في كيان}{الموارد البشرية \\ الموارد البشرية \\ المراسلات \\ مكتبة الملفات \\ المهام}
\stake{ما يتكرر في كل المؤسسات تجده جاهزاً. وما يخص مؤسستك وحدها هو ما تبنيه بأدوات البناء.}
\slink{main}{الدليل الرئيسي}{ar/index.html}
\end{rcslide}

\section{من كل ورقة إلى أداة}
\label{sec:paper-map}
\begin{rcslide}{من كل ورقة إلى أداة}
\sphase{البناء}
\slead{ما بقي بعد التطبيقات الجاهزة تبنيه، وكل ورقة لها أداة تناسبها.}
\scompare{الورقة}{دفتر تتابع فيه شيئاً \\ استمارة تُعبّأ \\ «وقّع وحوّل للمدير» \\ خدمة يطلبها المراجعون \\ خطاب يُكتب كل مرة \\ تقرير شهري}{الأداة في كيان}{كيان \\ نموذج \\ مسار موافقات أو إجراء داخلي \\ خدمة على البوابة \\ قالب مستند \\ لوحة معلومات}
\stake{ابدأ دائماً بالكيان: فالنموذج يكتب فيه، والإجراء يبدأ منه، واللوحة تحسب منه.}
\slink{builder}{دليل أدوات البناء}{ar/builder.html}
\end{rcslide}

\section{مثال كامل: الكشف على المنشآت}
\label{sec:paper-example}
\begin{rcslide}{مثال كامل: الكشف على المنشآت}
\sphase{مثال}
\slead{وزارة تكشف على المنشآت التجارية: يكلّف الباحث مفتشاً، فيعبّئ استمارة الكشف، ويقرر المدير.}
\schain{الفريق ← ملف المنشأة ← استمارة الكشف ← الموافقة ← النشر ← التكليف ← لوحة المؤشرات}
\scard{people}{الفريق}{الباحث والمفتش والمدير في الموارد البشرية.}
\scard{entity}{ملف المنشأة}{كيان يجمع كشوف كل منشأة.}
\scard{form}{استمارة الكشف}{نموذج يعبّئه المفتش، وعند اعتماده يتحدّث ملف المنشأة.}
\scard{chart}{لوحة المؤشرات}{يتابع فيها الباحث المنشآت والمخالفات.}
\stake{هذا هو الترتيب نفسه لأي إجراء: الناس، ثم السجل، ثم الاستمارة، ثم القرار، ثم الأرقام.}
\slink{inspection}{طبّقه خطوة خطوة: الكشف على المنشآت}{ar/index.html}
\end{rcslide}

\section{جرّب قبل أن تعمّم}
\label{sec:paper-pilot}
\begin{rcslide}{جرّب قبل أن تعمّم}
\sphase{التطبيق}
\slead{كل ما تبنيه يبقى مسودة حتى تنشره، فاستفد من ذلك.}
\schain{تبني ← تجرّب بسجلات حقيقية ← يجرّبه فريق صغير ← تصحّح ← تنشره للجميع}
\scard{people}{فريق صغير أولاً}{من يعرفون العمل ويصبرون على التعديل.}
\scard{paper}{سجلات حقيقية}{حالات فعلية من عملكم، لا أمثلة مصطنعة.}
\scard{flow}{الورق بالتوازي}{لفترة قصيرة، حتى تثق بالنتيجة.}
\stake{أخطاء التجربة رخيصة؛ أخطاء ما بعد التعميم مكلفة.}
\slink{builder}{المراجعة والنشر}{ar/builder-publish.html}
\end{rcslide}

\section{تمرين: حوّل إجراءً من عملك}
\label{sec:paper-exercise}
\begin{rcslide}{تمرين: حوّل إجراءً من عملك}
\sphase{التمرين}
\slead{اختر إجراءً ورقياً تعرفه جيداً، وأجب عن هذه الأسئلة الخمسة على ورقة.}
\scard{entity}{1. ما الذي تتابعه؟}{هذا هو الكيان: منشأة؟ طلب؟ مستفيد؟}
\scard{form}{2. ما الاستمارة؟}{أسئلتها، وما الإلزامي منها.}
\scard{people}{3. من يعبّئ، ومن يقرر؟}{المناصب لا الأسماء.}
\scard{flow}{4. ماذا يحدث بعد القرار؟}{مهمة؟ إشعار؟ مستند؟}
\scard{chart}{5. ما الرقم الذي يريده المدير؟}{هذه هي اللوحة.}
\stake{إجاباتك الخمس هي مخطط البناء كاملاً. قارنها بالمثال في الشريحة التالية.}
\end{rcslide}

\section{حل مقترح: طلب شراء}
\label{sec:paper-solution}
\begin{rcslide}{حل مقترح: طلب شراء}
\sphase{التمرين}
\slead{مثال محلول لإجراء تعرفه أغلب المؤسسات: طلب شراء يعتمده المدير ثم المالية.}
\smetric{1. ما الذي تتابعه؟}{دفتر الطلبات}{كيان «طلبات الشراء»}
\smetric{2. ما الاستمارة؟}{ورقة طلب الشراء}{نموذج: المادة والكمية والسبب}
\smetric{3. من يقرر؟}{توقيعان على الورقة}{مسار موافقات: المدير ثم المالية}
\smetric{4. بعد القرار؟}{يُبلَّغ مسؤول المشتريات}{إجراء داخلي ينشئ له مهمة}
\smetric{5. الرقم المطلوب؟}{جرد يدوي آخر الشهر}{لوحة: الطلبات المفتوحة والمتأخرة}
\stake{لاحظ أن الموظفين ومناصبهم لم تُبنَ هنا؛ فهم موجودون جاهزين في الموارد البشرية.}
\end{rcslide}

\section{أخطاء شائعة}
\label{sec:paper-mistakes}
\begin{rcslide}{أخطاء شائعة}
\sphase{الخلاصة}
\slead{أغلب مشكلات التحول لا تأتي من الأدوات، بل من طريقة البدء.}
\scompare{تجنّب}{نسخ الاستمارة الورقية حرفياً \\ البدء بالنموذج قبل الكيان \\ التفكير بالأشخاص بدل الأدوار \\ النشر للجميع دون تجربة}{افعل}{تحسين الاستمارة وأنت تنقلها \\ الكيان أولاً، ثم النموذج \\ صِف من يقرر بمنصبه أو دوره \\ فريق صغير أولاً، ثم التعميم}
\stake{النقل فرصة لتحسين الإجراء نفسه، لا لنسخ مشكلاته إلى شاشة.}
\end{rcslide}
